% -----------------------------------------------------------------------------
% Exact Ward--Takahashi identity from local U(1) covariance
% -----------------------------------------------------------------------------

局域 $U(1)$ 对称性不仅固定最小耦合顶角，也约束完整电子二点核与光子二点核。把精确一粒子不可约有效作用量记为 $\Gamma_{\rm 1PI}$，在经典背景四势 $A_\mu$ 中定义电子逆传播子
\begin{equation}
\boxed{
S_A^{-1}(x,y)
\equiv
\left.
\frac{\delta^2\Gamma_{\rm 1PI}}
{\delta\bar\psi(x)\,\delta\psi(y)}
\right|_{\psi=\bar\psi=0}.}
\label{eq:qed_background_inverse_propagator_definition}
\end{equation}
它是完整电子二点函数的逆核，并允许外部经典电磁背景任意强。采用
\[
A_\mu' = A_\mu-\partial_\mu\chi,
\qquad
U_\chi(x)=\exp\!\left(\frac{\ii q_{\mathrm e}\chi(x)}{\hbar}\right),
\]
局域协变性要求
\begin{equation}
\boxed{
S_{A-\partial\chi}^{-1}(x,y)
=U_\chi(x)S_A^{-1}(x,y)U_\chi^{-1}(y).}
\label{eq:qed_inverse_propagator_gauge_covariance}
\end{equation}
这个关系不依赖微扰展开。规范固定只使光子二次核可逆，不改变费米子端点的局域相位变换律。

剥去电荷因子后，精确一粒子不可约电磁顶角定义为
\begin{equation}
\boxed{
\Gamma^\mu(x,y;z)
\equiv
-\frac{1}{q_{\mathrm e}}
\left.
\frac{\delta S_A^{-1}(x,y)}{\delta A_\mu(z)}
\right|_{A=0}.}
\label{eq:qed_exact_vertex_functional_definition}
\end{equation}
树级时它退化为 $\Gamma^\mu=\gamma^\mu$。对\eqnrefs{eq:qed_inverse_propagator_gauge_covariance} 作无穷小规范变换，得到坐标空间 Ward--Takahashi 恒等式
\begin{equation}
\boxed{
\partial_\mu^z\Gamma^\mu(x,y;z)
=-\frac{\ii}{\hbar}
\bigl[\delta^{(4)}(z-x)-\delta^{(4)}(z-y)\bigr]
S^{-1}(x,y).}
\label{eq:qed_ward_takahashi_coordinate}
\end{equation}

\begin{proposition}[规范协变逆核的坐标 Ward 恒等式]
假设 1PI 逆传播核对背景规范变换协变，泛函微分存在，且正则化不产生矢量规范反常。则有以下结论（原式 \eqnrefs{eq:qed_ward_takahashi_coordinate}）：
\begin{equation*}
\boxed{
\partial_\mu^z\Gamma^\mu(x,y;z)
=-\frac{\ii}{\hbar}
\bigl[\delta^{(4)}(z-x)-\delta^{(4)}(z-y)\bigr]
S^{-1}(x,y).}
\end{equation*}
\end{proposition}
\begin{proof}
正文\eqnrefs{eq:qed_inverse_propagator_gauge_covariance} 给出背景场变化时精确电子逆传播子的端点协变性。对任意无穷小规范函数 $\chi$ 比较等式两侧的一阶变化，便能把这种端点相位约束转写成顶角的坐标空间散度关系。

左端对背景场作泛函展开，
\begin{align*}
 S_{-\partial\chi}^{-1}(x,y)
 &=S_0^{-1}(x,y)
 -\int\dd^4 z\,
 \left.\frac{\delta S_A^{-1}(x,y)}{\delta A_\mu(z)}\right|_0
 \partial_\mu\chi(z)+O(\chi^2)\\
 &=S_0^{-1}(x,y)
 +q_{\mathrm e}\int\dd^4 z\,\Gamma^\mu(x,y;z)\partial_\mu\chi(z)+O(\chi^2)\\
 &=S_0^{-1}(x,y)
 -q_{\mathrm e}\int\dd^4 z\,[\partial_\mu^z\Gamma^\mu(x,y;z)]\chi(z)+O(\chi^2).
\end{align*}
右端的两个端点相位展开为
\begin{align*}
 U_\chi(x)S_0^{-1}(x,y)U_\chi^{-1}(y)
 &=S_0^{-1}(x,y)
 +\frac{\ii q_{\mathrm e}}{\hbar}[\chi(x)-\chi(y)]S_0^{-1}(x,y)+O(\chi^2).
\end{align*}
两个端点值分别可写成
\begin{equation*}
 \chi(x)=\int\dd^4z\,\delta^{(4)}(z-x)\chi(z),
 \qquad
 \chi(y)=\int\dd^4z\,\delta^{(4)}(z-y)\chi(z).
\end{equation*}
所以右端一阶变化等于
\begin{equation*}
 \frac{\ii q_{\mathrm e}}{\hbar}\int\dd^4z\,
 \bigl[\delta^{(4)}(z-x)-\delta^{(4)}(z-y)\bigr]
 S_0^{-1}(x,y)\chi(z).
\end{equation*}
与左端
\begin{equation*}
 -q_{\mathrm e}\int\dd^4z\,
 [\partial_\mu^z\Gamma^\mu(x,y;z)]\chi(z)
\end{equation*}
相等。由于 $\chi(z)$ 是任意光滑测试函数，两个分布作用在任意测试函数上的积分都相等，只能有
\begin{equation*}
 \partial_\mu^z\Gamma^\mu(x,y;z)
 =-\frac{\ii}{\hbar}
 \bigl[\delta^{(4)}(z-x)-\delta^{(4)}(z-y)\bigr]
 S_0^{-1}(x,y),
\end{equation*}
即正文\eqnrefs{eq:qed_ward_takahashi_coordinate}。这里的两个 $\delta$ 接触项正是端点相位 $\chi(x)-\chi(y)$ 的分布表示：一个支撑在电子线的入射端点，另一个支撑在出射端点，二者符号相反。
\end{proof}

两个接触项分别对应光子插入点与入射、出射费米子端点重合；相对负号来自两端相反的规范相位。

令入射电子动量为 $p^\mu$、出射电子动量为 $p^\mu+q^\mu$，按全文 $\exp(-\ii p\cdot x/\hbar)$ 的 Fourier 约定，\eqnrefs{eq:qed_ward_takahashi_coordinate} 等价于
\begin{equation}
\boxed{
q_\mu\Gamma^\mu(p+q,p)
=S^{-1}(p+q)-S^{-1}(p).}
\label{eq:qed_ward_takahashi_momentum}
\end{equation}
这里两边都是精确一粒子不可约核。任何固定圈阶都必须把 $S^{-1}$ 与 $\Gamma^\mu$ 按同一阶截断，才能保持这一约束。

在 $q\to0$ 且 $S^{-1}(p)$ 对外动量可微的区域，\eqnrefs{eq:qed_ward_takahashi_momentum} 给出
\begin{equation}
\boxed{
\Gamma^\mu(p,p)
=\frac{\partial S^{-1}(p)}{\partial p_\mu}.}
\label{eq:qed_ward_zero_momentum_vertex_derivative}
\end{equation}

\begin{proposition}[零动量顶角与逆传播核导数]
设 Ward 恒等式成立，逆传播核在所论动量可微，零转移动量极限存在；奇异点需要保留分布边界值。则有以下结论（原式 \eqnrefs{eq:qed_ward_zero_momentum_vertex_derivative}）：
\begin{equation*}
\boxed{
\Gamma^\mu(p,p)
=\frac{\partial S^{-1}(p)}{\partial p_\mu}.}
\end{equation*}
\end{proposition}
\begin{proof}
按正文固定的 $\exp(-\ii p\cdot x/\hbar)$ 约定，把平移不变的逆传播子和三点顶角展开为
\begin{align*}
S^{-1}(x,y)
&=\int_p
 \ee^{-\ii p\cdot x/\hbar}
 \ee^{+\ii p\cdot y/\hbar}S^{-1}(p),\\
\Gamma^\mu(x,y;z)
&=\int_p\!\int_q
 \ee^{-\ii(p+q)\cdot x/\hbar}
 \ee^{+\ii p\cdot y/\hbar}
 \ee^{+\ii q\cdot z/\hbar}
 \Gamma^\mu(p+q,p),
\end{align*}
其中 $\int_p\equiv\int \dd^4 p/(2\pi\hbar)^4$。正文\eqnrefs{eq:qed_ward_takahashi_coordinate} 左边的 $z$ 导数只作用在最后一个 Fourier 相位上，因此
\begin{align*}
\partial_\mu^z\Gamma^\mu(x,y;z)
&=\int_p\!\int_q
 \ee^{-\ii(p+q)\cdot x/\hbar}
 \ee^{+\ii p\cdot y/\hbar}
 \ee^{+\ii q\cdot z/\hbar}
 \frac{\ii q_\mu}{\hbar}
 \Gamma^\mu(p+q,p).
\end{align*}
两个接触项分别使用
\begin{equation*}
\delta^{(4)}(z-x)=\int_q\ee^{+\ii q\cdot(z-x)/\hbar},
\qquad
\delta^{(4)}(z-y)=\int_q\ee^{+\ii q\cdot(z-y)/\hbar}.
\end{equation*}
第一项给
\begin{align*}
\delta^{(4)}(z-x)S^{-1}(x,y)
=\int_p\!\int_q
 \ee^{-\ii(p+q)\cdot x/\hbar}
 \ee^{+\ii p\cdot y/\hbar}
 \ee^{+\ii q\cdot z/\hbar}S^{-1}(p),
\end{align*}
而第二项把逆传播子的积分变量改写成 $p+q$ 后成为
\begin{align*}
\delta^{(4)}(z-y)S^{-1}(x,y)
=\int_p\!\int_q
 \ee^{-\ii(p+q)\cdot x/\hbar}
 \ee^{+\ii p\cdot y/\hbar}
 \ee^{+\ii q\cdot z/\hbar}S^{-1}(p+q).
\end{align*}
于是坐标空间恒等式两侧具有同一个 Fourier 核。比较系数，
\begin{equation*}
\frac{\ii}{\hbar}q_\mu\Gamma^\mu(p+q,p)
=-\frac{\ii}{\hbar}\left[S^{-1}(p)-S^{-1}(p+q)\right],
\end{equation*}
约去共同的 $\ii/\hbar$，得到正文\eqnrefs{eq:qed_ward_takahashi_momentum}：
\begin{equation*}
 q_\mu\Gamma^\mu(p+q,p)
 =S^{-1}(p+q)-S^{-1}(p).
\end{equation*}

再取 $q\to0$，并使用正文已经说明的可微条件。右边作 Taylor 展开：
\begin{equation*}
 S^{-1}(p+q)-S^{-1}(p)
 =q_\mu\frac{\partial S^{-1}(p)}{\partial p_\mu}+O(q^2).
\end{equation*}
左边的精确顶角在同一点展开为
\begin{equation*}
 q_\mu\Gamma^\mu(p+q,p)
 =q_\mu\Gamma^\mu(p,p)+O(q^2).
\end{equation*}
两边对任意足够小的四矢量 $q_\mu$ 都相等，因此比较每个独立的一阶系数得到
\begin{equation*}
 \Gamma^\mu(p,p)=\frac{\partial S^{-1}(p)}{\partial p_\mu},
\end{equation*}
即正文\eqnrefs{eq:qed_ward_zero_momentum_vertex_derivative}。
\end{proof}

因此电子波函数重整化与最小耦合顶角具有同一紫外乘法因子，
\begin{equation}
\boxed{Z_1=Z_2.}
\label{eq:qed_ward_z1_equals_z2}
\end{equation}

\begin{proposition}[零动量 Ward 恒等式的反项约束]
设顶角与波函数采用相同重整化条件，且所论矢量 Ward 恒等式无反常。则有以下结论（原式 \eqnrefs{eq:qed_ward_z1_equals_z2}）：
\begin{equation*}
\boxed{Z_1=Z_2.}
\end{equation*}
\end{proposition}
\begin{proof}
完整 1PI 顶角与完整电子逆传播子满足
\[
q_\mu\Gamma_B^\mu(p+q,p)
=
S_B^{-1}(p+q)-S_B^{-1}(p).
\]
令 \(q\to0\)，对 \(q_\mu\) 作一阶展开即得正文
\eqnrefs{eq:qed_ward_zero_momentum_vertex_derivative}：
\[
\Gamma_B^\mu(p,p)
=
\frac{\partial S_B^{-1}(p)}{\partial p_\mu}.
\]
这个关系作用在完整裸量上，因此任何紫外发散都必须在等式两边以相同方式出现。

要从这个恒等式抽取重整化常数，最干净的做法是只投影它的局域紫外发散部分，而不把方案依赖的有限项混进来。由表面发散度和 Lorentz 协变性，电子二点 1PI 核的局域紫外发散最多具有
\begin{equation*}
 \bigl[S_B^{-1}(p)\bigr]_{\rm UV,loc}
 =A_{\rm UV}\,\slashed p
 +B_{\rm UV}\,m_Rc,
\end{equation*}
其中 $A_{\rm UV},B_{\rm UV}$ 可以含 $1/\epsilon$ 极点，但不再依赖外动量；更高次的 $p$ 幂会对应不可重整化的局域算符。对 $p_\mu$ 求导后，质量型发散自动消失：
\begin{equation*}
 \frac{\partial}{\partial p_\mu}
 \bigl[S_B^{-1}(p)\bigr]_{\rm UV,loc}
 =A_{\rm UV}\gamma^\mu.
\end{equation*}

另一方面，零光子动量处的三点 1PI 顶角在 QED 中只有同维数的局域 Dirac 结构
\begin{equation*}
 \bigl[\Gamma_B^\mu(p,p)\bigr]_{\rm UV,loc}
 =C_{\rm UV}\gamma^\mu.
\end{equation*}
把这两个局域发散部分代入精确零动量 Ward--Takahashi 恒等式
\begin{equation*}
 \Gamma_B^\mu(p,p)
 =\frac{\partial S_B^{-1}(p)}{\partial p_\mu},
\end{equation*}
于是恒等式在唯一允许的局域基矢 $\gamma^\mu$ 上给出：
\begin{equation*}
 C_{\rm UV}\gamma^\mu=A_{\rm UV}\gamma^\mu
 \quad\Longrightarrow\quad
 C_{\rm UV}=A_{\rm UV}.
\end{equation*}

现在把反项写成
\begin{equation*}
 \delta_2\,\bar\psi_R\,\ii\hbar c\slashed\partial\psi_R,
 \qquad
 -\delta_1 q_Rc\,\bar\psi_R\gamma^\mu\psi_R A_{\mu,R},
 \qquad
 \delta_i\equiv Z_i-1.
\end{equation*}
重整化条件要求反项逐个抵消上述局域紫外系数，因此在同一符号约定下
\begin{equation*}
 \delta_2=-A_{\rm UV},
 \qquad
 \delta_1=-C_{\rm UV}.
\end{equation*}
由 $C_{\rm UV}=A_{\rm UV}$ 立即有
\begin{equation*}
 \delta_1=\delta_2,
 \qquad
 \boxed{Z_1=Z_2.}
\end{equation*}
质量反项不参与这一投影，因为 $\partial(m_Rc)/\partial p_\mu=0$；有限顶角与有限自能导数则继续由同一个 Ward--Takahashi 恒等式彼此约束。

因此，在保持局域 $U(1)$ Ward 恒等式的调节与重整化方案中，$Z_1=Z_2$ 是全阶约束。结合裸电荷与重整化电荷的关系，独立的电荷重整化信息只剩光子场因子 $Z_3$。
\end{proof}

物理电荷归一化进一步固定在壳 Dirac 形状因子 $F_1(0)=1$；Pauli 形状因子 $F_2(0)$ 不由纵向 Ward 恒等式决定，必须由顶角的横向有限部分计算。

同一个局域对称性也约束光子二点的一粒子不可约量子修正。把真空极化张量定义为有效作用量中的二次核
\begin{equation}
\boxed{
\Gamma_{\Pi}^{(2)}[A]
=-\frac12\int\frac{\dd^4 q}{(2\pi\hbar)^4}\,
A_\mu(-q)\Pi^{\mu\nu}(q)A_\nu(q).}
\label{eq:qed-vacpol-quadratic-kernel-dp68}
\end{equation}
这里自由光子核和规范固定项不包含在 \(\Pi^{\mu\nu}\) 中。量子修正本身必须对 \(A_\mu\to A_\mu-\partial_\mu\chi\) 不变，因此
\begin{equation}
\boxed{q_\mu\Pi^{\mu\nu}(q)=0.}
\label{eq:qed_vacuum_polarization_transversality_from_1pi}
\end{equation}

\begin{proposition}[规范不变二次核的横向性]
设规范不变的 1PI 有效作用量二阶可微，背景为平移不变真空，且边界项可消去。则有以下结论（原式 \eqnrefs{eq:qed_vacuum_polarization_transversality_from_1pi}）：
\begin{equation*}
\boxed{q_\mu\Pi^{\mu\nu}(q)=0.}
\end{equation*}
\end{proposition}
\begin{proof}
正文\eqnrefs{eq:qed-vacpol-quadratic-kernel-dp68} 把真空极化定义成光子有效作用量中的量子二次核。对它施加无穷小规范变换后，任意规范函数 \(\chi(q)\) 前的系数必须消失；这会直接限制 \(\Pi^{\mu\nu}\) 的纵向分量。

Fourier 空间的无穷小变换为
\begin{equation*}
\delta A_\mu(q)=-\frac{\ii}{\hbar}q_\mu\chi(q).
\end{equation*}
对正文\eqnrefs{eq:qed-vacpol-quadratic-kernel-dp68} 作一阶变分。二次核满足 $\Pi^{\mu\nu}(q)=\Pi^{\nu\mu}(-q)$；第二项作 $q\mapsto-q$ 后可与第一项合并，严格得到
\begin{equation*}
\delta\Gamma_{\Pi}^{(2)}
=-\frac{\ii}{\hbar}
\int\frac{\dd^4 q}{(2\pi\hbar)^4}\,
\chi(-q)q_\mu\Pi^{\mu\nu}(q)A_\nu(q).
\end{equation*}
规范函数 \(\chi\) 与测试场 \(A_\nu\) 可以独立选择，而量子有效作用量必须对任意 \(\chi\) 保持不变，因此只有
\begin{equation*}
q_\mu\Pi^{\mu\nu}(q)=0
\end{equation*}
才能使变分恒为零，这就是正文\eqnrefs{eq:qed_vacuum_polarization_transversality_from_1pi}。
\end{proof}

在 Lorentz 不变真空中，横向性把它唯一约化为
\begin{equation}
\boxed{
\Pi^{\mu\nu}(q)
=\left(q^2\eta^{\mu\nu}-q^\mu q^\nu\right)\Pi(q^2).}
\label{eq:qed_vacuum_polarization_transverse_decomposition}
\end{equation}
因此真空极化的张量结构、$Z_1=Z_2$ 与外光子纵向极化解耦都是同一局域 $U(1)$ 约束的不同投影；具体圈积分只计算对称性未固定的标量函数和有限横向结构。
